\documentclass[10pt,a4paper]{amsart}
\usepackage[margin=19mm]{geometry}
\usepackage{amsmath,amssymb,mathtools}
\usepackage[scr=rsfs,cal=euler]{mathalfa}
\usepackage{xfrac}
\usepackage[unicode,hidelinks,bookmarks=false]{hyperref}
\newcommand{\ie}{i.e.\ }
\newcommand{\cf}{cf.\ }
\newcommand{\resp}{respectively\ }
\newcommand{\Subsection}{\subsection}
\providecommand{\nobreakdash}{\nobreak\hbox{-}}
\setlength{\emergencystretch}{2em}
\multlinegap0pt
\usepackage{bm}
\usepackage{bbm}
\usepackage{dsfont}
\usepackage{xspace}
\usepackage{cancel}
\usepackage{mathtools}
\usepackage{amsthm}
\makeatletter
\@ifundefined{theorem}{\newtheorem{theorem}{Theorem}}{}
\makeatother
\title[Joint numerical radius of Tuples: Extreme points, subdiffere]{Joint numerical radius of Tuples: Extreme points, subdifferential set and Gateaux derivative}
\author{Arpita Mal}
\date{Source: arXiv:2507.04700v1 (CC BY 4.0)}
\begin{document}
\maketitle

\noindent\emph{Source and licence.} Joint numerical radius of Tuples: Extreme points, subdifferential set and Gateaux derivative. Version arXiv:2507.04700v1 (CC BY 4.0), distributed under the Creative Commons Attribution 4.0 International licence, as recorded in the arXiv licence metadata for this exact version and shown on the abstract page https://arxiv.org/abs/2507.04700v1. Full rights notice: licenses/arxiv-2507.04700-CC-BY-4.0.txt.\par\smallskip

\noindent\textit{Editorial scope.} Counted unit: the Theorem whose source label is \texttt{th-04} in the article (source0.tex lines 471--480, source label \texttt{th-04}), delivered together with the article's own complete original proof of it (source0.tex lines 481--512, one \texttt{proof} environment of 32 lines). No mathematics is written, repaired or re-derived: statement and proof are the article's own text line for line. Required context retained uncounted: none. Every definition, display and label of those lines is retained unchanged. Omitted from this package and not redistributed: 1--470, 513--611. Nothing inside a retained excerpt is omitted or altered: every retained body is delivered exactly as the article's lines are written, so no omission receipt of any kind accompanies this package. Citations: the retained excerpts carry no citation command at all, so no bibliography entry of the article is transcribed and no reference is redistributed. Reference adaptation: none was needed and none was made; every reference inside the delivered unit resolves inside it, so no unresolved reference or citation is produced. Printed slips kept literal: no printed word, symbol, hypothesis or number of the counted statement or of its proof was altered. Layout and class: the article's own class is \texttt{amsart} with packages including eucal; the wrapper is the lane's pinned \texttt{amsart} editorial wrapper at 10pt on A4, which restates the theorem-like environments the retained text needs and adds the article's own macro definitions that the retained text needs (none), with extra wrapper packages (bm, bbm, dsfont, xspace, cancel, mathtools). The delivered numbering is the unit's own. Assets: no retained span contains a figure environment, a graphics-inclusion command, a PDF-page inclusion, a table or a TikZ picture; no image file is copied into this package, the package declares no asset dependency and no image file is redistributed with it. Licence: version arXiv:2507.04700v1 of this article is distributed under the Creative Commons Attribution 4.0 International licence, as recorded in the arXiv licence metadata for this exact version and shown on the abstract page https://arxiv.org/abs/2507.04700v1. A full rights notice is delivered at licenses/arxiv-2507.04700-CC-BY-4.0.txt. Characters: every retained excerpt is pure ASCII, so the delivered engine, XeLaTeX with its default Latin Modern text font, reports no missing character.
\begin{theorem}\label{th-04}
		Let $\dim(\mathcal{X})<\infty.$  Suppose $\mathcal{T}=(T_1,\ldots,T_d)\in \mathcal{Z}$ and $w_p(\mathcal{T})\neq 0.$ Then the following are equivalent.\\
		\rm(i) The numerical radius $w_p(\cdot)$ of $\mathcal{Z}$ is Gateaux differentiable at $\mathcal{T}.$\\
		\rm(ii) There exists $(x_0,x_0^*)\in \Pi$ such that $x_0\in E_{\mathcal{X}}, x_0^*\in E_{\mathcal{X}^*}$ and $$M_{w_p(\mathcal{T})}=\{(\mu x_0, \overline{\mu} x_0^*):\mu\in \mathbb{F},|\mu|=1\}.$$
		Moreover, in this case, for each $\mathcal{S}=(S_1,\ldots,S_d) \in \mathcal{Z},$
	\begin{eqnarray*}
		G(\mathcal{T},\mathcal{S})
		&=&\frac{1}{w_p(\mathcal{T})^{p-1}}\sum_{i=1}^d\Re\bigg(\overline{x_0^*(T_ix_0)}|x_0^*(T_ix_0)|^{p-2}x_0^*(S_ix_0) \bigg).
	\end{eqnarray*}
\end{theorem}

\begin{proof}
	Without loss of generality, assume that $w_p(\mathcal{T})=1.$ \\
(i) $\Rightarrow $ (ii). Since $w_p(\cdot)$ is Gateaux differentiable at $\mathcal{T},$ $J_{w_p}(\mathcal{T})$ is singleton. Suppose $$J_{w_p}(\mathcal{T})=\bigg\{\bigg(\overline{x^*(T_1x)}|x^*(T_1x)|^{p-2}x^*\otimes x,\ldots,\overline{x^*(T_dx)}|x^*(T_dx)|^{p-2}x^*\otimes x\bigg)\bigg\}$$
for some $(x,x^*)\in M_{w_p(\mathcal{T})},$ where $x\in E_{\mathcal{X}}, x^*\in E_{\mathcal{X}^*}.$ We show that if
 $(y,y^*)\in M_{w_p(\mathcal{T})},$ then $(y,y^*)=(\mu x,\overline{\mu}x^*)$ for some scalar $\mu$ with $|\mu|=1.$ Clearly,
 \[\bigg(\overline{y^*(T_1y)}|y^*(T_1y)|^{p-2}y^*\otimes y,\ldots,\overline{y^*(T_dy)}|y^*(T_dy)|^{p-2}y^*\otimes y\bigg)\in J_{w_p}(\mathcal{T}).\]  
 Therefore for all $1\leq i\leq d$, 
 \begin{equation}\label{eq-01}
 \overline{y^*(T_iy)}|y^*(T_iy)|^{p-2}y^*\otimes y=\overline{x^*(T_ix)}|x^*(T_ix)|^{p-2}x^*\otimes x.
\end{equation}
 If possible, suppose that $\{x,y\}$ is linearly independent. Then there exists $S\in \mathcal{L}(\mathcal{X})$ such that $Sx=0$ and $Sy=y.$ Now from (\ref{eq-01}),
 \begin{eqnarray*}
 	 \overline{y^*(T_iy)}|y^*(T_iy)|^{p-2}y^*\otimes y(S)&=&\overline{x^*(T_ix)}|x^*(T_ix)|^{p-2}x^*\otimes x(S)\\
 	 \Rightarrow \overline{y^*(T_iy)}|y^*(T_iy)|^{p-2}&=&0, ~(\text{since } y^*(y)=1)\\
 	 \Rightarrow y^*(T_iy)&=&0,
 	\end{eqnarray*}
 which contradicts that $(y,y^*)\in M_{w_p(\mathcal{T})}.$ Therefore, assume that $y=\mu x.$ Since $y,x\in S_{\mathcal{X}},$ so $|\mu|=1.$ Now if possible, assume that $\{x^*,y^*\}$ is linearly independent. Choose $u\in \ker(x^*)\setminus \ker(y^*).$ Consider $S\in \mathcal{L}(\mathcal{X})$ defined by $S(z)=x^*(z)u$ for all $z\in \mathcal{X}.$ Again from (\ref{eq-01}), 
  \begin{eqnarray*}
 	\overline{y^*(T_iy)}|y^*(T_iy)|^{p-2}y^*\otimes y(S)&=&\overline{x^*(T_ix)}|x^*(T_ix)|^{p-2}x^*\otimes x(S)\\
 	\Rightarrow \overline{y^*(T_iy)}|y^*(T_iy)|^{p-2}y^*(\mu u)&=&0, ~(\text{since } x^*(x)=1)\\
 	\Rightarrow y^*(T_iy)&=&0,
 \end{eqnarray*}
 which contradicts that $(y,y^*)\in M_{w_p(\mathcal{T})}.$ Therefore, assume that $y^*=\lambda x^*.$  Now,
 \[1=y^*(y)=\lambda x^*(\mu x)=\lambda \mu\quad \Rightarrow \quad \lambda =\overline{\mu},\]
 which implies that $(y,y^*)=(\mu x,\overline{\mu}x^*).$ This proves (ii).\\
 
  (ii) $\Rightarrow$ (i).  In this case, clearly $J_{w_p}(\mathcal{T})$ is singleton. In particular,
  $$J_{w_p}(\mathcal{T})=\bigg\{\bigg(\overline{x_0^*(T_1x_0)}|x_0^*(T_1x_0)|^{p-2}x_0^*\otimes x_0,\ldots,\overline{x_0^*(T_dx_0)}|x_0^*(T_dx_0)|^{p-2}x_0^*\otimes x_0\bigg)\bigg\}$$
   and so $w_p(\cdot)$ is Gateaux differentiable at $\mathcal{T}.$\\
   
   The rest of the part follows from the fact that, $	G(\mathcal{T},\mathcal{S})=\Re(f(\mathcal{S})),$ if $J_{w_p}(\mathcal{T})=\{f\}.$ 
\end{proof}

\end{document}
